% Bagean: metodhe
% Bab: induksi
% Pérangan: relasi

\documentclass[../../../include/open-logic-section]{subfiles}

\begin{document}

\olfileid{mth}{ind}{rel}

\olsection{Relasi lan Fungsi}

Yèn sawijining himpunan obyèk, kaya bilangan asli
utawa term trep, wis didhéfinisikaké kanthi induktif,
kita uga bisa ndhéfinisikaké \emph{relasi ing}
obyèk mau kanthi induksi. Upamané, delengen
gagasan iki: term trep~$t_1$ iku subterm
saka term trep~$t_2$ yèn muncul minangka
pérangané. Gunakna simbol:
$t_1 \sqsubseteq t_2$. Mesthi saben term
trep dadi subtermé dhéwé:
$t \sqsubseteq t$. Relasi iki bisa
didhéfinisikaké kanthi induktif:

\begin{defn}
  Relasi yèn term trep~$t_1$
  dadi subterm saka~$t_2$,
  $t_1
  \sqsubseteq t_2$,
  didhéfinisikaké kanthi induksi
  ing~$t_2$ mangkéné:
  \begin{enumerate}
  \item Yèn $t_2$ aksara,
    $t_1 \sqsubseteq t_2$ yèn
    lan mung yèn $t_1 = t_2$.
  \item Yèn $t_2$ awujud
    $[s_1 \circ s_2]$, mula
    $t_1 \sqsubseteq t_2$ yèn
    lan mung yèn $t_1 =
    t_2$, $t_1 \sqsubseteq s_1$,
    utawa $t_1 \sqsubseteq s_2$.
  \end{enumerate}
\end{defn}

Upamané, dhéfinisi iki nuduhaké
yèn $\mathrm{a}
\sqsubseteq [\mathrm{b} \circ \mathrm{a}]$.
Miturut (2),
$\mathrm{a} \sqsubseteq [\mathrm{b} \circ \mathrm{a}]$
yèn lan mung yèn
$\mathrm{a}
= [\mathrm{b} \circ \mathrm{a}]$,
utawa $\mathrm{a} \sqsubseteq \mathrm{b}$,
utawa $\mathrm{a} \sqsubseteq \mathrm{a}$.
Rong kamungkinan kapisan salah:
$\mathrm{a}$ cetha ora padha karo
$[\mathrm{b} \circ
  \mathrm{a}]$, lan miturut~(1),
$\mathrm{a} \sqsubseteq \mathrm{b}$
yèn lan mung yèn
$\mathrm{a} = \mathrm{b}$,
kang uga salah. Nanging miturut~(1),
$\mathrm{a} \sqsubseteq \mathrm{a}$
yèn lan mung yèn
$\mathrm{a} = \mathrm{a}$,
kang bener.

Wigati dielingi yèn kasil oraé
dhéfinisi iki gumantung marang
kasunyatan kang durung dibuktèkaké:
saben term trep~$t$ iku aksara
tunggal, utawa ana term trep
$s_1$ lan $s_2$ kang
\emph{ditemtokaké kanthi siji cara}
saéngga $t = [s_1 \circ s_2]$.
``Ditemtokaké kanthi siji cara''
ateges yèn $t = [s_1 \circ s_2]$,
ora bisa \emph{uga}
$= [r_1 \circ r_2]$
kanthi $s_1 \neq r_1$
utawa $s_2 \neq r_2$.
Yèn bisa, klausa~(2)
mbokmenawa bakal nentang
awaké dhéwé: maca $t_2$
minangka $[s_1 \circ s_2]$
bisa mènèhi $t_1
\sqsubseteq t_2$, nanging
maca $t_2$ minangka
$[r_1 \circ r_2]$ bisa
mènèhi ora $t_1 \sqsubseteq t_2$.
Sadurungé mbuktekaké yèn iki
ora mungkin, ayo deleng conto
kang nuduhaké yèn bab iki
\emph{bisa} kelakon.

\begin{defn}
  Dhéfinisikna \emph{term tanpa kurung}
  kanthi induktif:
  \begin{enumerate}
  \item Saben aksara iku
    term tanpa kurung.
  \item Yèn $s_1$ lan $s_2$
    term tanpa kurung,
    $s_1 \circ s_2$ uga
    term tanpa kurung.
  \item Ora ana liya kang
    kalebu term tanpa kurung.
  \end{enumerate}
\end{defn}

Conto term tanpa kurung yaiku
$\mathrm{a}$,
$\mathrm{b} \circ
\mathrm{d}$, lan
$\mathrm{b} \circ \mathrm{a} \circ \mathrm{b}$.
Yèn ``subterm'' kanggo term
tanpa kurung didhéfinisikaké
kaya ing ndhuwur, klausa
kapindho dadi
\begin{center}
  Yèn $t_2 = s_1 \circ s_2$,
  mula $t_1 \sqsubseteq t_2$
  yèn lan mung yèn
  $t_1 = t_2$,
  $t_1 \sqsubseteq s_1$,
  utawa $t_1 \sqsubseteq s_2$.
\end{center}

Saiki $\mathrm{b} \circ \mathrm{a} \circ \mathrm{b}$
awujud $s_1
\circ s_2$ kanthi
\begin{align*}
  s_1 & = \mathrm{b} \text{ lan} & s_2 & = \mathrm{a} \circ \mathrm{b}.
\intertext{Iki uga awujud $r_1 \circ r_2$ kanthi}
  r_1 & = \mathrm{b} \circ \mathrm{a} \text{ lan} & r_2 &= \mathrm{b}.
\end{align*}
Apa $\mathrm{a} \circ \mathrm{b}$
iku subterm saka
$\mathrm{b} \circ
\mathrm{a} \circ \mathrm{b}$?
Miturut pamacaan kapisan,
wangsulané iya, nanging
miturut pamacaan kapindho
ora.

Sipat yèn cara mbangun
term trep saka term trep liya
iku tunggal diarani
\emph{mung bisa diwaca kanthi siji cara}.
Amarga dhéfinisi induktif
kanggo relasi ing obyèk
kaya iki wigati, kita kudu
mbuktekaké sipat mau.

\begin{prop}
  Anggepen $t$ term trep.
  Mula $t$ aksara tunggal,
  utawa ana term trep
  $s_1$, $s_2$ kang
  ditemtokaké kanthi siji
  cara saéngga~$t =
  [s_1 \circ s_2]$.
\end{prop}

\begin{proof}
  Yèn $t$ aksara tunggal,
  syarat iku kacukupan.
  Dadi anggepen $t$
  dudu aksara tunggal.
  Miturut dhéfinisi induktif,
  $t$ kudu awujud
  $[s_1 \circ s_2]$
  kanggo term trep
  $s_1$ lan~$s_2$.
  Kari mbuktekaké yèn
  loro-loroné ditemtokaké
  kanthi siji cara: yèn
  $t = [r_1 \circ r_2]$,
  mula $s_1 = r_1$
  lan $s_2 = r_2$.

  Anggepen $t = [s_1 \circ s_2]$
  lan uga $t = [r_1 \circ r_2]$
  kanggo term trep
  $s_1$, $s_2$, $r_1$, $r_2$.
  Kita kudu nuduhaké
  $s_1 = r_1$ lan
  $s_2 = r_2$. Kaping pisan,
  $s_1$ lan $r_1$
  mesthi padha, amarga
  yèn ora, salah sijiné
  dadi bagéan wiwitan sejati
  saka sijiné. Nanging
  miturut \olref[sti]{prop:initial},
  iki ora mungkin yèn
  $s_1$ lan~$r_1$ padha-padha
  term trep. Yèn
  $s_1 = r_1$, cetha
  $s_2 = r_2$ uga.
\end{proof}

Fungsi uga bisa didhéfinisikaké
kanthi induktif. Upamané,
fungsi~$f$ kang ngirim saben
term trep menyang tingkat
jeroné susunan $[\dots]$
kang paling gedhé bisa
didhéfinisikaké mangkéné:

\begin{defn}
  \ollabel{defn:depth}
  \emph{Jeroning} term trep,
  $f(t)$, didhéfinisikaké
  kanthi induktif:
  \[
  f(t) = \begin{cases}
    0 & \text{ yèn $t$ aksara}\\
    \max(f(s_1), f(s_2)) + 1 & \text{ yèn $t = [s_1 \circ s_2]$.}
  \end{cases}
  \]
\end{defn}

Upamané,
\begin{align*}
  f([\mathrm{a} \circ \mathrm{b}]) & =
    \max(f(\mathrm{a}),f(\mathrm{b})) + 1 = \\
  &= \max(0, 0) + 1 = 1, \text{ lan}\\
  f([[\mathrm{a} \circ \mathrm{b}] \circ \mathrm{c}]) & =
    \max(f([\mathrm{a} \circ \mathrm{b}]), f(\mathrm{c})) + 1 = \\
  & = \max(1,0) + 1 = 2.
\end{align*}

Ing kéné kita nganggep
$s_1$ lan $s_2$
term trep, lan migunakaké
kasunyatan yèn saben
term trep iku aksara
utawa awujud
$[s_1 \circ s_2]$.
Wigati maneh yèn wujud
iki mung ana siji cara.
Kanggo ngerti sebabé,
delengen maneh term tanpa
kurung. ``Dhéfinisi'' kang
cocog bakal awujud:
\[
  g(t) =
  \begin{cases}
    0 & \text{ yèn $t$ aksara}\\
   \max(g(s_1), g(s_2)) + 1 & \text{ yèn $t = s_1 \circ s_2$.}
  \end{cases}
\]
Saiki delengen term tanpa
kurung $\mathrm{a} \circ \mathrm{b} \circ
\mathrm{c} \circ \mathrm{d}$.
Term iki bisa diwaca
luwih saka siji cara,
upamané $s_1 \circ s_2$
kanthi
\begin{align*}
  s_1 & = \mathrm{a} \text{ lan} &
  s_2 & = \mathrm{b} \circ \mathrm{c} \circ \mathrm{d},
\intertext{utawa minangka $r_1 \circ r_2$ kanthi}
  r_1 & = \mathrm{a} \circ \mathrm{b} \text{ lan} &
  r_2 &= \mathrm{c} \circ \mathrm{d}.
\end{align*}
Yèn $g$ diétung manut
pamacaan kapisan, kita olèh
\begin{align*}
  g(s_1 \circ s_2) & =
  \max(g(\mathrm{a}), g(\mathrm{b} \circ \mathrm{c} \circ \mathrm{d})) + 1 =\\ & =
  \max(0,2) + 1 = 3
  \intertext{nanging manut pamacaan liya,}
  g(r_1 \circ r_2) & =
  \max(g(\mathrm{a} \circ \mathrm{b}), g(\mathrm{c} \circ \mathrm{d})) + 1=\\ &
  = \max(1,1) + 1 = 2
\end{align*}
Fungsi kudu tansah ngasilaké
nilai tunggal, mula
``dhéfinisi'' kanggo~$g$
iku ora ndhéfinisikaké
fungsi babar pisan.

\begin{prob}
  Wènèhana dhéfinisi
  induktif kanggo fungsi $l$,
  kanthi $l(t)$ cacah
  simbol ing term trep~$t$.
\end{prob}

\begin{prob}
  Buktèkna kanthi induksi
  struktural ing term trep
  $t$ yèn $f(t) < l(t)$,
  kanthi $l(t)$ cacah simbol
  ing~$t$ lan $f(t)$
  jeroning $t$ kaya
  didhéfinisikaké ing
  \olref[mth][ind][rel]{defn:depth}.
\end{prob}

\end{document}
